\documentclass[11pt,a4paper]{article}
\usepackage[margin=2.4cm]{geometry}
\usepackage{amsmath,amssymb}
\usepackage{booktabs}
\usepackage{microtype}
\usepackage[hidelinks]{hyperref}
\usepackage{parskip}
\setlength{\parskip}{0.55em}
\title{\textbf{PANO Phase 0 Sample Document}}
\author{PANO Project}
\date{2026}
\begin{document}
\maketitle

\begin{abstract}
This sample PDF is used to exercise continuous scrolling, text selection, highlights,
geometric annotations, Markdown notes and offline math rendering in the PANO Phase 0 demo.
The source PDF is read-only; all annotations are stored beside it in a \texttt{.pano} directory.
\end{abstract}

\section{Reading and Annotation}
PANO treats the PDF as an immutable source. A highlight is not a rectangle painted over a
screenshot: it keeps the selected text range and every selection polygon reported by the PDF
engine. That makes the annotation meaningful when the document is reopened at another zoom
level or with a different display.

\subsection{Text marks}
Select a sentence with the text tool and apply a highlight, underline or strike-out. The mark
inherits the text range and the page coordinates in PDF points. Moving the window or rotating
the view must not move the mark relative to the glyphs.

\subsection{Geometric marks}
Lines, rectangles, ellipses, pencil paths and the highlighter are editable objects. A single
drag creates one undoable edit. The eraser deletes the topmost object in one step, and
\texttt{Ctrl+Z} restores it.

\begin{equation}
  E = mc^2
\end{equation}

\begin{equation}
  \int_a^b f(x)\,\mathrm{d}x \approx \sum_{i=1}^{n} f(x_i)\,\Delta x
\end{equation}

\section{Notes Anchored to the Source}
A note is attached to an annotation instead of replacing the sentence it discusses. The note
panel stores Markdown, and the preview supports headings, lists, tables, code spans, local
images and math. Inline math such as $m\ddot{u}+ku=0$ and display math such as

\[
  \nabla \cdot \vec{E} = \frac{\rho}{\varepsilon_0}
\]

remain fully offline.

\subsection{A small table}
\begin{table}[h]
\centering
\begin{tabular}{lll}
\toprule
Quantity & Symbol & Unit \\
\midrule
Mass & $m$ & kg \\
Stiffness & $k$ & N/m \\
Damping & $c$ & N$\cdot$s/m \\
\bottomrule
\end{tabular}
\caption{Parameters used by the example.}
\end{table}

\section{Storage and Recovery}
The annotation file is JSON and is committed with an atomic replace. A SHA-256 digest binds the
data directory to the exact PDF bytes. If the digest does not match, PANO stops instead of
silently rebinding the annotations to a different document.

Every drag, stroke or erase is one command on the undo stack. Note edits are restorable as well.
The undo history is deliberately not persisted: reopening the book gives a clean editing
session while keeping the saved annotations.

\section{Validation Checklist}
\begin{enumerate}
  \item Open this PDF and check that pages scroll continuously.
  \item Zoom with \texttt{Ctrl+wheel}; the point under the cursor should stay under it.
  \item Select text and create a highlight; rotate the view by 90 degrees.
  \item Draw a rectangle, change its colour and opacity, then undo the change.
  \item Attach a Markdown note, capture a page region as an image, and save.
  \item Close and reopen the PDF; the marks and notes must return.
\end{enumerate}

\newpage
\section{Long-form Reading Page}
This page exists to make continuous scrolling and page jumping observable. The quick brown fox
jumps over the lazy dog. The quick brown fox jumps over the lazy dog. The quick brown fox jumps
over the lazy dog. The quick brown fox jumps over the lazy dog. The quick brown fox jumps over
the lazy dog. The quick brown fox jumps over the lazy dog. The quick brown fox jumps over the
lazy dog. The quick brown fox jumps over the lazy dog.

\subsection{Selection across lines}
Text selection is deliberately limited to one page in Phase 0. A selection that begins near the
top of this paragraph and ends near its bottom should still produce one annotation whose
polygons cover every selected line. Starting a drag on one page and releasing on another page is
ignored rather than producing a corrupt cross-page object.

\subsection{Mathematical text}
The renderer used by notes is a local, deliberately small TeX subset. It is not a replacement for
a full typesetting engine; it exists so the core reading experience does not depend on a browser
engine or a network connection. Fractions, roots, Greek letters, sums, integrals, accents and
common operators are enough for engineering notes.

\begin{align}
  \frac{\partial u}{\partial t} &= \alpha \nabla^2 u, \\
  \sqrt{x^2+y^2} &\le \left(\frac{a+b}{c}\right)^2.
\end{align}

\section{Closing}
Use the note tool to place a pin in the margin, then write a Markdown note with a table, a code
snippet and a captured region. Save the document, close it, and open it again. The pin should
still be at the same page position, and the note should still point at the same source object.

\begin{quote}
An annotation should explain the source, not replace it.
\end{quote}

\end{document}
